% GENERATED FILE. Edit canonical lessons, then run npm run generate:books.
% canonical-source-hash: fnv1a64:db7b5c66d6836ff3
% canonical-lessons: GU-C548-tim, GU-C548-kaptan, GU-C548-mukablo, GU-C548-vijeta, GU-C548-jitadvu

\chapter{If We Win the Match}
\label{ch:gu-if-we-win-the-match}
\hlchaptermodality{548}

\begin{chapteropening}
\textbf{By the end of this chapter:} \emph{I can talk about a team, its captain and the winner.}

Complete the last lesson of chapter 548: I can talk about a team, its captain and the winner.
\end{chapteropening}

\section[\texorpdfstring{\gu{ટીમ} (ṭīm)}{ટીમ (ṭīm)}]{\gu{ટીમ} (ṭīm) — a team}

\label{lesson:GU-C548-tim}



\begin{quote}
Before the new one: say the Gujarati for rather, then the Gujarati for adverse.
\end{quote}

\subsection*{You'll want to know: \gu{ટીમ}}
\textbf{\gu{ટીમ}} — \emph{ṭīm} — \textquotedblleft{}a team\textquotedblright{}.

\begin{quote}
\gu{જો} \gu{અમારી} \gu{ટીમ} \gu{આજે} \gu{જીતશે} \gu{તો} \gu{અમે} \gu{મીઠાઈ} \gu{વહેંચીશું}.
\end{quote}

\begin{quote}
\emph{jo amārī ṭīm āje jītśe to ame mīṭhāī vahẽcīśũ.}
\end{quote}

\begin{quote}
If our team wins today, we'll hand out sweets.
\end{quote}

A feminine noun.

\begin{cousinweb}
\textbf{\gu{ટીમ}} (\emph{ṭīm}) is borrowed from English \emph{team}.
\end{cousinweb}

\begin{grammarlens}[title={What you've built}]
One more word for saying what happens if.
\end{grammarlens}

\subsection*{Your turn}
\begin{itemize}
\raggedright
  \item \emph{Say it:} \emph{ṭīm}
  \item \emph{Say it:} \emph{ṭīm}, once more
  \item \emph{Say it:} \emph{jo amārī ṭīm āje jītśe to ame mīṭhāī vahẽcīśũ.} — If our team wins today, we'll hand out sweets.
  \item \emph{Recall:} say the Gujarati for rather, then the Gujarati for adverse, then say \emph{ṭīm} again
\end{itemize}

\begin{tcolorbox}[breakable,colback=teal!4,colframe=teal!35!black,title={Before you move on}]
What does \gu{ટીમ} mean? (\textquotedblleft{}A team\textquotedblright{}.) Say it once more.
\end{tcolorbox}
\section[\texorpdfstring{\gu{કપ્તાન} (kaptān)}{કપ્તાન (kaptān)}]{\gu{કપ્તાન} (kaptān) — a captain}

\label{lesson:GU-C548-kaptan}



\begin{quote}
Before the new one: say the Gujarati for adverse, then the Gujarati for a team.
\end{quote}

\subsection*{You'll want to know: \gu{કપ્તાન}}
\textbf{\gu{કપ્તાન}} — \emph{kaptān} — \textquotedblleft{}a captain\textquotedblright{}.

\begin{quote}
\gu{જો} \gu{કપ્તાન} \gu{સારું} \gu{રમશે} \gu{તો} \gu{ટીમ} \gu{જીતશે}.
\end{quote}

\begin{quote}
\emph{jo kaptān sārũ ramśe to ṭīm jītśe.}
\end{quote}

\begin{quote}
If the captain plays well, the team will win.
\end{quote}

A masculine noun.

\begin{cousinweb}
\textbf{\gu{કપ્તાન}} (\emph{kaptān}) is usually traced to Portuguese \emph{capitão}, \textquotedblleft{}a captain\textquotedblright{}, the word of the ships that reached the Indian coast.
\end{cousinweb}

\begin{grammarlens}[title={What you've built}]
One more word for saying what happens if.
\end{grammarlens}

\subsection*{Your turn}
\begin{itemize}
\raggedright
  \item \emph{Say it:} \emph{kaptān}
  \item \emph{Say it:} \emph{kaptān}, once more
  \item \emph{Say it:} \emph{jo kaptān sārũ ramśe to ṭīm jītśe.} — If the captain plays well, the team will win.
  \item \emph{Recall:} say the Gujarati for adverse, then the Gujarati for a team, then say \emph{kaptān} again
\end{itemize}

\begin{tcolorbox}[breakable,colback=teal!4,colframe=teal!35!black,title={Before you move on}]
What does \gu{કપ્તાન} mean? (\textquotedblleft{}A captain\textquotedblright{}.) Say it once more.
\end{tcolorbox}
\section[\texorpdfstring{\gu{મુકાબલો} (mukāblo)}{મુકાબલો (mukāblo)}]{\gu{મુકાબલો} (mukāblo) — a match-up}

\label{lesson:GU-C548-mukablo}



\begin{quote}
Before the new one: say the Gujarati for a team, then the Gujarati for a captain.
\end{quote}

\subsection*{You'll want to know: \gu{મુકાબલો}}
\textbf{\gu{મુકાબલો}} — \emph{mukāblo} — \textquotedblleft{}a match-up\textquotedblright{}.

\begin{quote}
\gu{જો} \gu{કાલે} \gu{વરસાદ} \gu{પડશે} \gu{તો} \gu{મુકાબલો} \gu{રવિવારે} \gu{થશે}.
\end{quote}

\begin{quote}
\emph{jo kāle varsād paḍśe to mukāblo ravivāre thaśe.}
\end{quote}

\begin{quote}
If it rains tomorrow, the contest will be held on Sunday.
\end{quote}

A masculine noun.

\begin{cousinweb}
\textbf{\gu{મુકાબલો}} (\emph{mukāblo}) is usually traced to Arabic \emph{muqābala}, \textquotedblleft{}facing one another, an encounter\textquotedblright{}, through Persian.
\end{cousinweb}

\begin{grammarlens}[title={What you've built}]
One more word for saying what happens if.
\end{grammarlens}

\subsection*{Your turn}
\begin{itemize}
\raggedright
  \item \emph{Say it:} \emph{mukāblo}
  \item \emph{Say it:} \emph{mukāblo}, once more
  \item \emph{Say it:} \emph{jo kāle varsād paḍśe to mukāblo ravivāre thaśe.} — If it rains tomorrow, the contest will be held on Sunday.
  \item \emph{Recall:} say the Gujarati for a team, then the Gujarati for a captain, then say \emph{mukāblo} again
\end{itemize}

\begin{tcolorbox}[breakable,colback=teal!4,colframe=teal!35!black,title={Before you move on}]
What does \gu{મુકાબલો} mean? (\textquotedblleft{}A match-up\textquotedblright{}.) Say it once more.
\end{tcolorbox}
\section[\texorpdfstring{\gu{વિજેતા} (vijetā)}{વિજેતા (vijetā)}]{\gu{વિજેતા} (vijetā) — a winner}

\label{lesson:GU-C548-vijeta}



\begin{quote}
Before the new one: say the Gujarati for a captain, then the Gujarati for a match-up.
\end{quote}

\subsection*{You'll want to know: \gu{વિજેતા}}
\textbf{\gu{વિજેતા}} — \emph{vijetā} — \textquotedblleft{}a winner\textquotedblright{}.

\begin{quote}
\gu{જો} \gu{તું} \gu{છેલ્લી} \gu{બાજી} \gu{જીતીશ} \gu{તો} \gu{તું} \gu{વિજેતા} \gu{બનીશ}.
\end{quote}

\begin{quote}
\emph{jo tũ chellī bājī jītīś to tũ vijetā banīś.}
\end{quote}

\begin{quote}
If you win the last round, you'll be the winner.
\end{quote}

A masculine or feminine noun, by the person meant.

\begin{cousinweb}
\textbf{\gu{વિજેતા}} (\emph{vijetā}) is a tatsama, from Sanskrit \emph{vijetṛ}, \textquotedblleft{}a conqueror, a victor\textquotedblright{} (Monier-Williams).
\end{cousinweb}

\begin{grammarlens}[title={What you've built}]
One more word for saying what happens if.
\end{grammarlens}

\subsection*{Your turn}
\begin{itemize}
\raggedright
  \item \emph{Say it:} \emph{vijetā}
  \item \emph{Say it:} \emph{vijetā}, once more
  \item \emph{Say it:} \emph{jo tũ chellī bājī jītīś to tũ vijetā banīś.} — If you win the last round, you'll be the winner.
  \item \emph{Recall:} say the Gujarati for a captain, then the Gujarati for a match-up, then say \emph{vijetā} again
\end{itemize}

\begin{tcolorbox}[breakable,colback=teal!4,colframe=teal!35!black,title={Before you move on}]
What does \gu{વિજેતા} mean? (\textquotedblleft{}A winner\textquotedblright{}.) Say it once more.
\end{tcolorbox}
\section[\texorpdfstring{\gu{જિતાડવું} (jitāḍvũ)}{જિતાડવું (jitāḍvũ)}]{\gu{જિતાડવું} (jitāḍvũ) — to win it for (a team)}

\label{lesson:GU-C548-jitadvu}



\begin{quote}
Before the new one: say the Gujarati for a match-up, then the Gujarati for a winner.
\end{quote}

\subsection*{You'll want to know: \gu{જિતાડવું}}
\textbf{\gu{જિતાડવું}} — \emph{jitāḍvũ} — \textquotedblleft{}to win it for (a team)\textquotedblright{}.

\begin{quote}
\gu{જો} \gu{આપણો} \gu{કપ્તાન} \gu{સારું} \gu{રમત} \gu{તો} \gu{તે} \gu{ટીમને} \gu{જિતાડત}.
\end{quote}

\begin{quote}
\emph{jo āpṇo kaptān sārũ ramat to te ṭīmne jitāḍat.}
\end{quote}

\begin{quote}
If our captain had played well, he would have won it for the team.
\end{quote}

Past, with \gu{તેણે} (he or she …): \gu{તેણે} \gu{ટીમને} \gu{જિતાડી}, \emph{teṇe ṭīmne jitāḍī}.

\begin{cousinweb}
\textbf{\gu{જિતાડવું}} (\emph{jitāḍvũ}) is built from \textbf{\gu{જીતવું}}, \textquotedblleft{}to win\textquotedblright{}, with the causative \emph{-āḍ-} and a shortened vowel: to make someone win. CDIAL files \textbf{\gu{જીતવું}} under Sanskrit \emph{jita}, \textquotedblleft{}won\textquotedblright{} (CDIAL 5224).
\end{cousinweb}

\begin{grammarlens}[title={What you've built}]
One more word for saying what happens if.
\end{grammarlens}

\subsection*{Your turn}
\begin{itemize}
\raggedright
  \item \emph{Say it:} \emph{jitāḍvũ}
  \item \emph{Say it:} \emph{jitāḍvũ}, once more
  \item \emph{Say it:} \emph{jo āpṇo kaptān sārũ ramat to te ṭīmne jitāḍat.} — If our captain had played well, he would have won it for the team.
  \item \emph{Recall:} say the Gujarati for a match-up, then the Gujarati for a winner, then say \emph{jitāḍvũ} again
\end{itemize}

\begin{tcolorbox}[breakable,colback=teal!4,colframe=teal!35!black,title={Before you move on}]
What does \gu{જિતાડવું} mean? (\textquotedblleft{}To win it for (a team)\textquotedblright{}.) Say it once more.
\end{tcolorbox}
